\section{Conception mécanique}

\subsection{Architecture mécanique}

Le scanner se présente sous la forme d'une boîte fermée de dimensions
extérieures inférieures à $500\times500\times500$\,mm pour une masse
totale inférieure à 15\,kg (exigence \textbf{E21}). L'architecture
est pensée pour une intégration directe dans l'environnement du
MakeLab de la HEIG-VD (FC3).

Le système s'organise autour de trois sous-ensembles :

\begin{itemize}
  \item \textbf{Plateau rotatif motorisé} (FT7, \textbf{E15}) :
        support de l'objet à numériser, entraîné par un moteur
        pas-à-pas NEMA\,17 couplé à un étage de transmission
        (courroie ou réduction directe, à figer).
  \item \textbf{Bâti de mesure} : supporte la caméra Pi\,Camera\,3
        et le module laser ligne vert dans une géométrie fixe et
        calibrée (angle connu, exigence \textbf{E12}).
  \item \textbf{Carter} : couvercle amovible équipé d'un
        micro-interrupteur de sécurité (SW1) assurant la coupure
        automatique du laser à l'ouverture (\textbf{E19}).
\end{itemize}

\emph{À compléter : vue éclatée CAO, rendu d'assemblage et liste
des pièces mécaniques (BOM mécanique).}

\subsection{Zone de scan et périmètre fonctionnel}

La zone utile est dimensionnée pour accueillir un objet maximal de
$150\times150\times150$\,mm (\textbf{E3}). Conformément au périmètre
fonctionnel \textbf{E23}, le système cible les pièces simples :

\begin{itemize}
  \item sans cavité fermée (non accessible par la ligne de vue
        caméra/laser) ;
  \item sans contre-dépouille supérieure à $45^\circ$ par rapport
        à l'axe de rotation.
\end{itemize}

Cette limitation découle du principe de triangulation laser
mono-vue et doit être communiquée à l'utilisateur via l'interface.

\subsection{Choix des matériaux}

Les matériaux accessibles à l'utilisateur sont non toxiques et les
arêtes chanfreinées ou ébavurées (\textbf{E17}). Le carter est
envisagé en tôle pliée peinte ou panneau MDF, choix résultant d'un
compromis coût~/~rigidité~/~facilité d'usinage au MakeLab. Les
éléments structurels internes combinent profilés aluminium (type
$20\times20$) et pièces imprimées en PETG.

\emph{À compléter : fiche matériaux définitive après revue mécanique,
essais de rigidité et tolérances d'alignement.}

\subsection{Dimensionnement}

Le dimensionnement du plateau rotatif vise un couple suffisant pour
supporter un objet de masse pouvant atteindre 2\,kg (hypothèse
haute) à vitesse lente (inférieure à 1\,tour/min). Le couple
résistant comprend principalement le frottement des roulements et
un éventuel balourd. Le moteur NEMA\,17 sélectionné est
surdimensionné pour cette charge et offre une marge confortable
avec le microstepping 1/32 du driver DM320T.

\emph{À compléter : bilan des forces, vérification moteur + réduction,
dimensionnement de l'axe et des roulements.}

\subsection{Sécurité mécanique}

Deux mesures principales assurent la protection de l'utilisateur :

\begin{itemize}
  \item \textbf{Confinement des pièces mobiles} (\textbf{E18}) :
        plateau et moteur enfermés dans le carter fermé, aucune
        pièce en rotation n'est accessible pendant un scan.
  \item \textbf{Interlock de sécurité laser} (\textbf{E19}) :
        micro-interrupteur SW1 sur le couvercle, câblé de telle
        sorte que son ouverture coupe logiciellement la GPIO
        \emph{Laser on} via un passage en état \texttt{ERROR} de
        la machine d'états.
\end{itemize}

\subsection{Points ouverts}

\begin{itemize}
  \item Schéma CAO détaillé et plans de fabrication.
  \item Fixation du laser et réglage d'angle (calibration initiale).
  \item Matériau final du plateau et traitement de surface
        anti-réflexion (limitation des parasites lumineux, FT1).
  \item Tolérances d'alignement caméra/laser requises pour atteindre
        la précision \textbf{E24} ($\leq 2$\,mm).
\end{itemize}
